\documentclass[10pt,a4paper]{amsart}
\usepackage[margin=19mm]{geometry}
\usepackage{amsmath,amssymb,mathtools}
\usepackage[scr=rsfs,cal=euler]{mathalfa}
\usepackage{xfrac}
\usepackage[unicode,hidelinks,bookmarks=false]{hyperref}
\newcommand{\ie}{i.e.\ }
\newcommand{\cf}{cf.\ }
\newcommand{\resp}{respectively\ }
\newcommand{\Subsection}{\subsection}
\providecommand{\nobreakdash}{\nobreak\hbox{-}}
\setlength{\emergencystretch}{2em}
\multlinegap0pt
\usepackage{multirow}
\usepackage{amssymb}
\usepackage{xcolor}
\usepackage{booktabs}
\usepackage{algorithm}\usepackage{algpseudocode}
\usepackage{tikz}
\newtheorem{theorem}{Theorem}
\DeclareMathOperator*{\argmax}{argmax}
\title{Sequential Apportionments from Stationary Divisor Methods}
\author{Michael A. Jones, Brittany Ohlinger,, Jennifer Wilson}
\date{Source: arXiv:2512.20686v2 (CC BY 4.0)}
\begin{document}
\maketitle

\noindent\emph{Source and licence.} Sequential Apportionments from Stationary Divisor Methods. Version arXiv:2512.20686v2 (CC BY 4.0), distributed by arXiv under the Creative Commons Attribution 4.0 International licence, http://creativecommons.org/licenses/by/4.0/, as recorded in the arXiv licence metadata for this exact version and shown on the abstract page https://arxiv.org/abs/2512.20686v2. Full licence notice: \texttt{licenses/arxiv-2512.20686-CC-BY-4.0.txt}.\par\smallskip

\noindent\textit{Editorial scope.} Counted unit: the article's theorem jeff\_adams (source0.tex lines 181--182). The article's complete original proof of it is delivered with it (source0.tex lines 184--196, one \texttt{proof} environment of 13 lines). No mathematics is written, repaired or re-derived: the statement and the proof are the article's own lines, in the article's own notation, and no retained line is substituted or re-flowed.  No further source-local context is needed: every cross-reference of every retained line resolves inside the two delivered spans.  Nothing was omitted from any retained span, so the package carries no omission record.  No retained line contains a citation, so the wrapper carries no bibliography and no bibliography entry.  No retained line contains a figure environment, an image-inclusion command or drawing code, so no image is delivered and the package declares no asset dependency.  Editorial formatting only, in addition: an AMS \texttt{amsart} preamble that restates the article's own declarations and macros used by the retained lines, namely the environments \texttt{theorem} on separate counters, together with the standard packages those lines need.  The retained environments print in this document as Theorem 1 in order of appearance, whereas the article prints the same items with the numbering of its own full document.  The complete native source of the article as distributed in the arXiv e-print for this exact version is bundled unchanged as \texttt{sources/191590/source0.tex}.
\begin{theorem}\label{jeff_adams}
For every positive vote vector $\mathbf p = (p_1, p_2, \ldots, p_n)$, if $S_1(\mathbf p) = (s_1, s_2, \dots)$ is the sequential apportionment under D'Hondt's method, then the sequential apportionment under Adams' method is $S_0(\mathbf p) = (t_1, t_2, \dots) = (1, 2, \dots, n, s_1, s_2, s_3, \dots)$.  \end{theorem}

\begin{proof} 
We proceed by induction. First we show our base case that $t_{n+1} = s_1$.  Under Adams' method, each of the parties $1, 2,  \dots, n$ receive, in this order, one of the first $n$ seats, so that $t_i = i$ for $i = 1$ to $n$. Seat $n+1$ is given to the party maximizing $\frac{p_i}{1}$, which is party 1 because $p_1 \ge p_i$ for $i=2$ to $n$; hence, $t_{n+1} = 1$. Under D'Hondt's method, $s_1 = 1$ because $i=1$ maximizes the same ratio: $\frac{p_i}{1}$. And, $t_{n+1} = s_1$.

Assume that $t_{i+n} = s_i$ for $i = 2$ to $k$, then we show that $t_{k+n+1} = s_{k+1}$. Let $\mathbf a$ be the apportionment under the D'Hondt method for $h=k$. Then, $a_j$ is the number of $s_i = j$ for $i = 1$ to $k$.  By definition, seat $k+1$ goes to the party with the smallest index $i$ among those indices that maximize $\frac{p_i}{a_i+1}$, or equivalently,
\begin{equation*}
\min_i \argmax_{i = 1 \text{ to }n} \left\{ \frac{p_i}{a_i + 1} \right\}.
\end{equation*}
By the inductive hypothesis, $b_j = a_j + 1$ is party $j$'s apportionment under Adams' method for $h = k+n$. It follows that seat $k+n+1$ goes to party 
\begin{equation*}
\min_i \argmax_{i = 1 \text{ to }n} \left\{ \frac{p_i}{b_i} \right\} = \min_i \argmax_{i = 1 \text{ to }n} \left\{ \frac{p_i}{a_i + 1} \right\}.
\end{equation*}
Hence, $t_{k+n+1} = s_{k+1}$, completing the proof.
\end{proof}

\end{document}
